%%%%%%%%%%%%%%%%%%%%%%%%%%%%%%%%%%%%%%%%%%%%%%%%%%%%%%%%%%%%%%%%%%%%%%%
%%%% Список сокращений и условных обозначений                     %%%%
%%%%%%%%%%%%%%%%%%%%%%%%%%%%%%%%%%%%%%%%%%%%%%%%%%%%%%%%%%%%%%%%%%%%%%%

\noindent\textit{Список сокращений}

\begin{longtable}{@{}p{0.28\textwidth}p{0.66\textwidth}@{}}
БПЛА    & беспилотный летательный аппарат\\
ВАК     & Высшая аттестационная комиссия при Министерстве науки и~высшего образования Российской Федерации\\
ИПС     & информационно-поисковая система\\
ИПУ~РАН & Институт проблем управления им.\,В.\,А.\,Трапезникова Российской академии наук\\
ИТИС    & Институт информационных технологий и~интеллектуальных систем\\
КА      & космический аппарат\\
МКПУ    & Мультиконференция по~проблемам управления\\
НИЦ     & Научно-исследовательский центр\\
ОЗУ     & оперативное запоминающее устройство\\
ПЭВМ    & персональная электронно-вычислительная машина\\
РАН     & Российская академия наук\\
ROI     & область интереса (region of interest)\\
УБС     & Управление большими системами (серия конференций)\\
ФГБУ    & федеральное государственное бюджетное учреждение\\
ФГБОУ~ВО & федеральное государственное бюджетное образовательное учреждение высшего образования\\
ЭВМ     & электронно-вычислительная машина\\
ArcFace & функция потерь Additive Angular Margin Loss для распознавания лиц~/ метрического обучения\\
AUC     & площадь под ROC-кривой (Area Under Curve)\\
BCE     & бинарная кросс-энтропия (Binary Cross-Entropy)\\
C2f     & модуль Cross-Stage Partial bottleneck with two convolutions, fast (архитектурный блок YOLOv8)\\
CIoU    & функция потерь Complete Intersection over Union\\
CNN     & сверточная нейронная сеть (Convolutional Neural Network)\\
CUDA    & программно-аппаратная платформа параллельных вычислений NVIDIA\\
CUHK03  & набор данных Chinese University of Hong Kong (версия~03) для задачи реидентификации\\
DFL     & Distribution Focal Loss (функция потерь распределённой фокальной регрессии)\\
DSPA    & конференция <<Цифровая обработка сигналов и~её применение>> (Digital Signal Processing and its Applications)\\
EER     & уровень равной ошибки (Equal Error Rate)\\
FAR     & уровень ложного допуска (False Acceptance Rate)\\
FRR     & уровень ложного отклонения (False Rejection Rate)\\
GCN     & графовая сверточная сеть (Graph Convolutional Network)\\
GPU     & графический процессор (Graphics Processing Unit)\\
HOG     & гистограмма направленных градиентов (Histogram of Oriented Gradients)\\
ICCT    & International Conference on Information, Control, and Communication Technologies\\
IoU     & пересечение по~объединению ограничивающих рамок (Intersection over Union)\\
LDA     & линейный дискриминантный анализ (Linear Discriminant Analysis)\\
LFW     & набор данных Labeled Faces in the Wild\\
mAP     & средняя средняя точность (mean Average Precision)\\
MAP     & максимум апостериорной вероятности (Maximum A~Posteriori)\\
Market-1501 & набор данных для реидентификации людей\\
MARS    & набор данных Motion Analysis and Re-identification Set\\
MLSD    & International Conference on Management of Large-Scale System Development\\
NMS     & подавление немаксимумов (Non-Maximum Suppression)\\
PCA     & анализ главных компонент (Principal Component Analysis)\\
PR      & кривая <<точность~--- полнота>> (Precision-Recall)\\
ReID    & реидентификация (re-identification)\\
ROC     & характеристическая кривая (Receiver Operating Characteristic)\\
SOTA    & state-of-the-art (наилучший на~текущий момент результат)\\
SPPF    & ускоренный модуль пирамидального пулинга (Spatial Pyramid Pooling Fast)\\
SVM     & метод опорных векторов (Support Vector Machine)\\
TAR     & уровень истинного допуска (True Acceptance Rate)\\
t-SNE   & метод нелинейного снижения размерности (t-Distributed Stochastic Neighbor Embedding)\\
ViT     & архитектура Vision Transformer\\
YOLO~/ YOLOv8 & архитектура детектирования объектов (You Only Look Once) и~её 8-я версия\\
\end{longtable}

\vspace{1em}

\noindent\textit{Условные обозначения}

\begin{longtable}{@{}p{0.28\textwidth}p{0.66\textwidth}@{}}
$K$ & число компонентов многокомпонентной структуры\\
$C_k$ & $k$-й компонент структуры, $k=1,\ldots,K$ (в~частном случае~$K=2$: $C_1$~--- лицо, $C_2$~--- силуэт)\\
$\Phi_k$ & нейросетевой кодировщик (encoder) $k$-го компонента, формирующий векторное представление\\
$\mathbf{z}_k\in\mathbb{R}^{d_k}$ & вектор признакового представления (эмбеддинг) $k$-го компонента в~пространстве размерности~$d_k$\\
$\mathbf{f}$, $\mathbf{s}$ & в~рассматриваемом случае~$K=2$: эмбеддинги лица и~силуэта соответственно\\
$d_k$, $d_f$, $d_s$ & размерности признаковых пространств соответствующих компонентов ($d_f=128$, $d_s=256$)\\
$\mathcal{M}(t)\subseteq\{1,\ldots,K\}$ & множество компонентов, наблюдаемых в~момент времени~$t$\\
$Y$ & множество классов идентификации (галерея ИПС)\\
$y$, $y^{*}$ & класс идентификации и~его оценка по~правилу принятия решения\\
$P(y)$ & априорная вероятность класса~$y$\\
$P(\mathbf{z}_k\mid y)$ & правдоподобие наблюдения эмбеддинга~$\mathbf{z}_k$ при условии класса~$y$\\
$P(y\mid\{\mathbf{z}_k\})$ & апостериорная вероятность класса~$y$ по~наблюдаемым эмбеддингам\\
$\boldsymbol{\mu}_{y,k}$ & центроид (среднее) распределения эмбеддингов класса~$y$ в~признаковом пространстве компонента~$k$\\
$\Sigma_{y,k}$ & ковариационная матрица распределения эмбеддингов класса~$y$ в~признаковом пространстве компонента~$k$\\
$\mathcal{I}=\{I_i\}_{i=1}^{N_I}$ & множество входных изображений (кадров) видеопотока\\
$B_h$, $B_f$ & множества ограничивающих рамок силуэтов и~лиц на~изображении\\
$\mathcal{L}_0$ & базовая (исходная) функция потерь YOLOv8\\
$\mathcal{L}_{\text{inner}}$ & добавочный компонент функции потерь, штрафующий ошибки локализации вложенного компонента (в~случае ReID~--- лица)\\
$\mathbb{I}_{\text{inner}}(i)$ & индикаторная функция: $1$, если $i$-я ячейка содержит вложенный компонент, $0$~--- иначе\\
$\lambda_{\text{inner}}$ & управляющий коэффициент при компоненте~$\mathcal{L}_{\text{inner}}$ в~суммарной функции потерь\\
$\mathcal{L}_{\text{CIoU}}$, $\mathcal{L}_{\text{BCE}}$, $\mathcal{L}_{\text{DFL}}$ & компоненты стандартной функции потерь YOLOv8\\
$\mathcal{L}_{\text{ArcFace}}$ & функция потерь ArcFace для метрического обучения эмбеддингов\\
$N_{\text{pred}}$ & общее число предсказаний детектора, соответствующих объектам\\
$s$, $m$ & масштабный коэффициент и~аддитивная угловая margin-поправка в~ArcFace\\
\end{longtable}
